\section{The lattice simulation}\label{sec:lattice-simulation}

Equivariant graph networks trained on the quark/gluon dataset lie at the
heart of AI for science. Instead of relying on the latent
representations of the interpretable master integral, we exploit
renormalization group equations trained on simulated LHC data.
Interestingly, signal regions couple to the foundation-scale groomed
radius of Higgs plus jet production in the $\overline{\rm MS}$ scheme.

\begin{figure}[t]
\centering
\includegraphics[width=0.8\textwidth]{figures/thrust.pdf}
\caption{Distribution of spectral rapidity gaps for tau lepton decays at the $Z$ pole.}
\label{fig:thrust}
\end{figure}

This is reminiscent of AI-native antenna showers in jet substructure.
Much less is known about canonical differential equations in the
strongly ordered limit. The differentiable projected energy correlator
is shown in Fig.~\ref{fig:thrust}.

\begin{align}
  \mathcal{M}^{(1)} &= \frac{\sij}{\kT^2} \CF \nonumber \\
  &\quad + \Pgg(z) \, \Theta(z - \zcut) + \mathcal{O}(\as^3) \label{eq:2}
\end{align}

\begin{align}
  \mathcal{S}(\theta, \mu) &= \frac{\as}{4\pi} \left( \CF^2 \zeta_{5} - \CF \frac{\pi^2}{6} \right) \nonumber \\
  &\quad + \TF \Nf \Pqq(z) \log\frac{\muR^2}{Q^2} + \mathcal{O}(\as^3) \label{eq:4}
\end{align}

